\begingroup

\section{Complete all-degree principal-angle proof and literal tensor theta primitive}


\begingroup
\sloppy
\small

\noindent\textbf{Source role.} Complete written proof, with original source bytes and reversible notation/publication mapping retained.\par
\noindent\textbf{Exact revision.} \nolinkurl{b0634428d69f80dedd4f68766ef9e6468c390eaac86dfe58b08c21c8fb896e03}\par
\noindent\textbf{Complete source.} \path{sources/next_edition/all_degree_review_conversions_v1/sources/all_degree_angle_review/tau_all_degree_angle_independent_review_20260913.md}\par
\noindent\textbf{Source SHA-256.} \nolinkurl{b0634428d69f80dedd4f68766ef9e6468c390eaac86dfe58b08c21c8fb896e03}\par\medskip
\subsubsection{All-degree principal angles of the original theta quotient sections}\label{endpoint-all_degree_angle_review--all-degree-principal-angles-of-the-original-theta-quotient-sections}

13 September 2026. This is a complete independent derivation of the nonlinear arithmetic determinant comparison in every quotient degree. All spaces are the original polynomial spaces, all projections use the actual common interpolated source Gram, and every change of representative is exhibited as a theta boundary. The earlier sealed AT and AW sources are unchanged.

\subsubsection{1. Fixed original objects and the canonical projection relation}\label{endpoint-all_degree_angle_review--fixed-original-objects-and-the-canonical-projection-relation}

Retain a full actual finite zero packet \(h\) for \(g=2\xi\), an integer \(k\ge1\), its original tensor-sum annihilator \(\chi=\chi_{h,k}\) of degree \(q\ge1\), and

\[
E=\mathbb C[S]/(\chi),\quad H=\mathcal P_{2q},\quad
\mathcal P_N=\mathbb C[S]_{\le N},\quad\mathcal P_{-1}=0.
\tag{AA1}
\]

Use the original increasing-power remainder frame of \(E\) and the original source coordinate \(S=k/2+iu\). The two source densities are exactly

\[
\begin{aligned}
m_1(u)&=w_h^{*k}(u),&w_h(u)&=\frac{|(g/h)(1/2+iu)|^2}{2\pi},\\
m_0(u)&=\frac{(\sqrt{2\pi})^k}{c_{k/4}}
 \frac{|\Gamma(k/4+iu/2)|^2}{2\pi},&c_a&=2^{1-2a}\Gamma(2a).
\end{aligned}
\tag{AA2}
\]

Their literal masses \(\mu_h^k\) and \((2\pi)^{k/2}\) remain in all Grams. Put \(m_t=(1-t)m_0+tm_1\), \(0\le t\le1\), and let \(M(t)\) be its common Gram on \(H\). The source proofs establish positivity almost everywhere and every polynomial moment for the two densities. Therefore \(M(t)\) and every restricted Gram below are positive definite. Inner products are conjugate-linear in their first argument, with \(\langle x,y\rangle_t=x^*M(t)y\).

For \(q-1\le N\le2q\), retain the exact sequence

\[
0\to\mathcal P_{N-q}\xrightarrow{B_N=\times\chi}
\mathcal P_N\xrightarrow{J_N}E\to0.
\tag{AA3}
\]

Let \(P_N\) be the common-metric orthogonal projection onto \(\mathcal P_N\), and \(Q_N\) the projection onto \(\mathcal R_N=\chi\mathcal P_{N-q}\). In particular \(Q_{q-1}=0\). Since \(\mathcal R_N\subseteq\mathcal P_N\), the two projections commute and \(P_NQ_N=Q_NP_N=Q_N\). Hence

\[
N_N:=P_N-Q_N
\quad\text{is the orthogonal projection onto}\quad
K_N=\mathcal P_N\cap\mathcal R_N^\perp.
\tag{AA4}
\]

Polynomial division gives \(\dim K_N=q\). The original minimum section \(R_N:E\to\mathcal P_N\subset H\) satisfies

\[
J_NR_N=I_E,\quad\operatorname{im}R_N=K_N,
\quad G_N=R_N^*M(t)R_N>0,\quad V_N=\det G_N.
\tag{AA5}
\]

These objects retain their dependence on \(t\). We suppress it only inside pointwise equations.

For any two cutoffs \(q-1\le i\le j\le2q\), both sections have remainder equal to their input. Thus \(R_i-R_j\) has image in \(\mathcal R_j\). Meanwhile \(R_j\) lies in \(K_j\), orthogonal to that image. Orthogonal projection of \(R_i\) onto \(K_j\) consequently gives the exact relation

\[
R_j=N_jR_i,\qquad
R_i^*MR_j=G_j,
\qquad
G_i-G_j=(R_i-R_j)^*M(R_i-R_j)\succeq0.
\tag{AA6}
\]

No equality between the different canonical sections is assumed. The same formula is valid at \(i=j\), with zero difference.

Let \(H_j=B_j^*M B_j\), with every map understood through its original coefficient inclusion in \(H\). The relation coefficient and its full energy are

\[
\begin{aligned}
X_{ij}&=H_j^{-1}B_j^*M R_i:E\to\mathcal P_{j-q},\\
R_i-R_j&=B_jX_{ij},\qquad
G_i-G_j=X_{ij}^*H_jX_{ij}.
\end{aligned}
\tag{AA7}
\]

Indeed \(Q_jR_i=B_jH_j^{-1}B_j^*MR_i\), and the splitting in \(\mathcal P_j\) shows that this is \(R_i-R_j\). If the relation domain is zero, the formula uses its unique zero map and no inverse on a nonzero space. This is the actual source relation energy that will become the squared sine data.

\subsubsection{2. Finite principal-angle calculation with all Gram factors}\label{endpoint-all_degree_angle_review--finite-principal-angle-calculation-with-all-gram-factors}

All square roots in this section are the positive Hermitian square roots in the fixed coefficient frame. Define the isometric frames

\[
F_i=R_iG_i^{-1/2}:\mathbb C^q\to(H,M),\qquad
F_j=R_jG_j^{-1/2}:\mathbb C^q\to(H,M).
\tag{AA8}
\]

They satisfy \(F_i^*MF_i=F_j^*MF_j=I\); this equality retains the original Grams rather than assigning new source norms. Their exact cross matrix is

\[
Z_{ij}=F_i^*MF_j
=G_i^{-1/2}G_jG_j^{-1/2}
=G_i^{-1/2}G_j^{1/2}.
\tag{AA9}
\]

The order of its factors matters. Multiplication by its actual adjoint gives

\[
A_{ij}:=Z_{ij}Z_{ij}^*
=G_i^{-1/2}G_jG_i^{-1/2},
\qquad 0<A_{ij}\preceq I.
\tag{AA10}
\]

Its positivity follows from \(G_j>0\), and its upper bound from AA6. Also \(I-A_{ij}=G_i^{-1/2}X_{ij}^*H_jX_{ij}G_i^{-1/2}\), retaining the complete source boundary energy.

Let \(v_1,\ldots,v_q\) be an orthonormal eigenbasis of \(A_{ij}\), with eigenvalues \(c_1^2,\ldots,c_q^2\) in \((0,1]\), including every multiplicity. Define

\[
x_a=F_iv_a\in K_i,\qquad
y_a=c_a^{-1}N_jx_a\in K_j,
\qquad 1\le a\le q.
\tag{AA11}
\]

There is no division by zero because of AA10. Since \(F_i^*MN_jF_i=Z_{ij}Z_{ij}^*=A_{ij}\), direct calculation gives

\[
\langle x_a,x_b\rangle_t=\delta_{ab},\quad
\langle y_a,y_b\rangle_t=\delta_{ab},\quad
\langle x_a,y_b\rangle_t=c_b\delta_{ab}.
\tag{AA12}
\]

For example the second equality follows from \(\langle N_jx_a,N_jx_b\rangle_t
=\langle x_a,N_jx_b\rangle_t=c_b^2\delta_{ab}\), divided by \(c_ac_b\). The \(y_a\) form a basis of \(K_j\) because they are \(q\) orthonormal vectors in that space.

Define \(\theta_a\in[0,\pi/2)\) by \(c_a=\cos\theta_a\). These are the complete principal angles of the two actual section spaces, as witnessed by the paired frames AA11--AA12. Their squared cosines are exactly the eigenvalues of AA10, with no interchange of noncommuting square roots.

If \(c_a=1\), then the orthogonal projection \(N_jx_a\) has the full norm of \(x_a\), so \(x_a\in K_j\) and \(y_a=x_a\). Conversely every \(x\in K_i\cap K_j\) is fixed by \(N_iN_j\) on \(K_i\), and its coefficient vector is an eigenvector of AA10 of eigenvalue one. Hence

\[
\#\{a:\theta_a=0\}=\dim(K_i\cap K_j),
\tag{AA13}
\]

including all repeated eigenvalues of one.

For every \(a\) with \(s_a=\sin\theta_a>0\), put \(z_a=(y_a-c_ax_a)/s_a\). Equations AA12 imply that the vectors \(z_a\) are orthonormal, orthogonal to every \(x_b\), and that different two-dimensional spaces \(\operatorname{span}(x_a,z_a)\) are orthogonal. They are also orthogonal to every common vector with \(c_b=1\). On each such plane, the difference of the original projections is exactly

\[
(N_i-N_j)|_{\operatorname{span}(x_a,z_a)}
=\begin{pmatrix}s_a^2&-s_ac_a\\-s_ac_a&-s_a^2\end{pmatrix}.
\tag{AA14}
\]

This follows because \(N_i\) is projection onto \(x_a\) on that plane, while \(N_j\) is projection onto \(y_a=c_ax_a+s_az_a\). The trace is zero and the determinant is \(-s_a^2\), giving eigenvalues \(+s_a,-s_a\). On the common vectors and on \((K_i+K_j)^\perp\), the difference is zero. Thus if \(r=\#\{a:\theta_a>0\}\), its full spectrum on \(H\) is

\[
\operatorname{spec}(N_i-N_j)
=\{+\sin\theta_a,-\sin\theta_a:\theta_a>0\}
 \sqcup\{0^{[\dim H-2r]}\}.
\tag{AA15}
\]

This proves the complete finite angle-to-projection spectral relation, including zero angles, repeated positive angles, and the remaining zero eigenvalues.

Finally, taking determinants of AA10 gives

\[
\prod_{a=1}^q\cos^2\theta_a
=\frac{V_j}{V_i},\qquad
\log\frac{V_i}{V_j}
=\sum_{a=1}^q -\log\cos^2\theta_a\ge0.
\tag{AA16}
\]

Every determinant is computed in the unchanged remainder frame. In particular this identity retains the full original quotient metric, not merely its eigenvalue locations.

\subsubsection{3. The two original cutoff pairs and their forced common direction}\label{endpoint-all_degree_angle_review--the-two-original-cutoff-pairs-and-their-forced-common-direction}

Use the original grouped pairs

\[
(i_0,j_1)=(q-1,2q),\qquad
(i_1,j_0)=(q,2q-1).
\tag{AA17}
\]

The second pair is ordered for \(q\ge2\); for \(q=1\) its two indices both equal one. Apply AA11--AA16 to both pairs, retaining all \(2q\) angles, and define the original endpoint quantity

\[
\begin{aligned}
\mathcal B(t)
&=\log\frac{V_{q-1}(t)V_q(t)}{V_{2q-1}(t)V_{2q}(t)}\\
&=\log\frac{V_{q-1}(t)}{V_{2q}(t)}
 +\log\frac{V_q(t)}{V_{2q-1}(t)}
=\sum_{a=1}^{2q}x_a(t),\quad
x_a=-\log\cos^2\theta_a\ge0.
\end{aligned}
\tag{AA18}
\]

At least one angle in the second pair is exactly zero. To prove this from the original spaces, both \(K_q\) and \(K_{2q-1}\) lie in

\[
L=\mathcal P_{2q-1}\cap(\chi\mathcal P_0)^\perp,
\qquad\dim L=2q-1.
\tag{AA19}
\]

They lie in that hyperplane because each is orthogonal to its own relation space, which contains \(\chi\mathcal P_0\). The latter is a nonzero line inside \(\mathcal P_{2q-1}\), so it cuts that \(2q\)-dimensional space by exactly one dimension in a positive form. Since each \(K\) has dimension \(q\), \(\dim(K_q\cap K_{2q-1})=2q-\dim(K_q+K_{2q-1})\ge1\). Now AA13 proves the claim. At \(q=1\), the two one-dimensional spaces are identical, in agreement with their identical cutoffs.

Keep the complete angle list. At each \(t\), remove one of these forced entries \(x_a=0\) only for the purpose of writing the following finite sum, and put

\[
m=2q-1.
\tag{AA20}
\]

The remaining \(m\) entries include any further zero angles and still sum to \(\mathcal B(t)\). No continuous choice of the removed entry or of individual eigenvectors is needed: the next inequality is a pointwise symmetric inequality in this entire list. The smooth quantity differentiated below is the original log determinant, not individual angle branches.

\subsubsection{4. Direct trace bound and finite concavity calculation}\label{endpoint-all_degree_angle_review--direct-trace-bound-and-finite-concavity-calculation}

Define the unchanged relative operators

\[
D=M(0)^{-1}M(1),\quad
C(t)=M(t)^{-1}(M(1)-M(0))
=((1-t)I+tD)^{-1}(D-I).
\tag{AA21}
\]

The last identity follows by factoring \(M(t)=M(0)((1-t)I+tD)\) in its actual order. The operator \(D\) is positive self-adjoint for \(M(0)\), and \(C(t)\) is self-adjoint for \(M(t)\) because \(M(t)C(t)=M(1)-M(0)\) is Hermitian. Let its extreme eigenvalues be \(c_{\min}(t),c_{\max}(t)\), and set \(w(t)=c_{\max}(t)-c_{\min}(t)\ge0\).

For any unit vector \(v\), expansion in an orthonormal eigenbasis of \(C(t)\) gives \(\langle v,C(t)v\rangle_t=\sum c_j|v_j|^2\), between those two extremes. Apply this fact to the positive and negative eigenvectors in AA15. Each pair contributes at most \(\sin\theta_a\,w(t)\) in absolute value. Summing and using the triangle inequality proves directly

\[
|\operatorname{Tr}((N_i-N_j)C(t))|
\le\left(\sum_{a=1}^q\sin\theta_a\right)w(t).
\tag{AA22}
\]

No external trace extremum theorem, or commutation with the section projections, is used.

The original determinant variation is \(\frac d{dt}\log V_N=\operatorname{Tr}(N_N C(t))\). Its four original signs give exactly

\[
\dot{\mathcal B}
=\operatorname{Tr}((N_{q-1}-N_{2q}+N_q-N_{2q-1})C(t)).
\tag{AA23}
\]

It follows either by differentiating the finite determinant ratio from the exact sequence, or from the retained AT12. Apply AA22 to the two pairs in AA17. Using AA18 and AA20 yields the full angle-profile bound

\[
|\dot{\mathcal B}(t)|
\le\left(\sum_{a=1}^{2q}\sin\theta_a(t)\right)w(t)
=\left(\sum_{a=1}^{m}\sqrt{1-e^{-x_a(t)}}\right)w(t).
\tag{AA24}
\]

For \(x>0\), direct differentiation of \(f(x)=\sqrt{1-e^{-x}}\) gives

\[
f'(x)=\frac{e^{-x}}{2\sqrt{1-e^{-x}}},\qquad
f''(x)=-\frac{e^{-x}(2-e^{-x})}{4(1-e^{-x})^{3/2}}<0.
\tag{AA25}
\]

Thus its derivative decreases, so its secant slopes decrease, proving concavity on \((0,\infty)\). Iterating the two-point concavity inequality proves the finite equal-weight inequality \(m^{-1}\sum f(y_a)\le f(m^{-1}\sum y_a)\) for positive \(y_a\). For nonnegative \(x_a\), apply it to \(y_a=x_a+\varepsilon\) and let \(\varepsilon\downarrow0\); continuity at zero, with \(f(0)=0\), gives the same result. Therefore AA24 implies

\[
\boxed{
|\dot{\mathcal B}(t)|
\le m\sqrt{1-e^{-\mathcal B(t)/m}}\,w(t),
\qquad m=2q-1.
}
\tag{AA26}
\]

At \(m=1\) the finite concavity step is an equality. Additional zero angles in higher degree are still present in AA24; AA26 is the stated consequence using only their sum and the guaranteed one zero angle.

\subsubsection{5. Strict positivity from the actual vertical-line moment source}\label{endpoint-all_degree_angle_review--strict-positivity-from-the-actual-vertical-line-moment-source}

Every term in AA18 is nonnegative. Suppose that \(\mathcal B(t)=0\) at one value of \(t\). Then every angle of the first pair is zero, so AA13 makes its two \(q\)-dimensional spaces equal:

\[
K_{q-1}=K_{2q}.
\tag{AA27}
\]

The first is \(\mathcal P_{q-1}\), because its relation domain is zero. It contains the original constant polynomial one. The second is orthogonal to \(\chi\mathcal P_q\). Thus \(1\perp\chi\mathcal P_q\).

Write the original polynomial in full as \(\chi(S)=\sum_{j=0}^q a_jS^j\), \(a_q=1\), and define its actual reflected-conjugate polynomial

\[
\chi_k^{\#}(S)=\sum_{j=0}^q\overline{a_j}(k-S)^j.
\tag{AA28}
\]

Its degree is \(q\), with leading coefficient \((-1)^q\overline{a_q}\); this sign is retained. On the unchanged integration line, \(\overline S=k-S\), and so

\[
\chi_k^{\#}(S)=\overline{\chi(S)},\qquad
\chi(S)\chi_k^{\#}(S)=|\chi(S)|^2.
\tag{AA29}
\]

The polynomial product belongs to the literal relation subspace \(\chi\mathcal P_q\). Orthogonality of one to that space would give

\[
0=\langle1,\chi\chi_k^{\#}\rangle_t
=\int_{\mathbb R}|\chi(k/2+iu)|^2m_t(u)\,du>0,
\tag{AA30}
\]

a contradiction. The strict inequality is the positive original Gram of the nonzero polynomial \(\chi\), and all integrals exist by the retained moments. This argument uses the actual source measure and vertical-line multiplication identity. It is not a claim about every abstract positive coefficient metric. We have proved

\[
\mathcal B(t)>0\qquad(0\le t\le1).
\tag{AA31}
\]

The finite positive Grams and their inverses are smooth near the compact path, so \(\mathcal B\) is smooth and attains a strictly positive minimum on it.

\subsubsection{6. Nonlinear endpoint transport in every quotient degree}\label{endpoint-all_degree_angle_review--nonlinear-endpoint-transport-in-every-quotient-degree}

Let the full ordered positive spectrum of the original relative Gram be

\[
0<b_1\le\cdots\le b_{2q+1},\qquad
\kappa=b_{2q+1}/b_1.
\tag{AA32}
\]

From AA21 its eigenvalues are \(c_j(t)=(b_j-1)/(1-t+tb_j)\). Differentiating with respect to \(b_j\) gives \((1-t+tb_j)^{-2}>0\), so the same extreme indices apply for every \(t\), including repeated values and values equal to one. Since each \(c_j\) is the derivative of \(\log(1-t+tb_j)\),

\[
\int_0^1 w(t)dt=\log b_{2q+1}-\log b_1=\log\kappa.
\tag{AA33}
\]

For the fixed integer \(m=2q-1\), define on \(B>0\)

\[
\Phi_m(B)=2\operatorname{arcosh}(e^{B/(2m)}),\qquad
\Phi_m'(B)=\frac1{m\sqrt{1-e^{-B/m}}}.
\tag{AA34}
\]

The derivative follows by differentiating \(\cosh(\Phi_m(B)/2)=e^{B/(2m)}\), using the nonnegative real inverse branch of \(\cosh\). The square root divisor is strictly positive on the original path by AA31. Multiply AA26 by AA34 and integrate the absolute derivative. Equation AA33 gives the all-degree nonlinear comparison

\[
\boxed{
\left|2\operatorname{arcosh}(e^{\mathcal B(1)/(2m)})
      -2\operatorname{arcosh}(e^{\mathcal B(0)/(2m)})\right|
\le\log\kappa,
\qquad m=2q-1.
}
\tag{AA35}
\]

The original values \(\mathcal B(0)=\mathcal B^{\Gamma}_{h,k}\) and \(\mathcal B(1)=\mathcal B^{\mathrm{ar}}_{h,k}\) are kept. No monotonicity of \(\mathcal B(t)\), individual angles, or individual eigenvectors was assumed.

An equivalent endpoint interval retains both branches explicitly. Set

\[
A_0=\operatorname{arcosh}(e^{\mathcal B(0)/(2m)}),\qquad
L=\tfrac12\log\kappa.
\tag{AA36}
\]

The endpoint inverse-cosh value lies in \([\max(0,A_0-L),A_0+L]\). Since \(A\mapsto2m\log\cosh A\) has derivative \(2m\tanh A\ge0\) for \(A\ge0\),

\[
\boxed{
2m\log\cosh(\max\{0,A_0-L\})
\le\mathcal B(1)\le2m\log\cosh(A_0+L).
}
\tag{AA37}
\]

The lower bound may be zero while the actual endpoint is strictly positive. When \(q=1\), \(m=1\), the second pair has its one zero angle and AA35 is the previously proved one-dimensional formula. For any higher \(q\), including the actual quartet dimension, the same proof applies. It establishes no uniform-in-\(k\) bound for \(\kappa\).

The new estimate implies the earlier bound \(|\mathcal B(1)-\mathcal B(0)|\le m\log\kappa\): the derivative of the inverse function \(B=2m\log\cosh(\Phi/2)\) is \(m\tanh(\Phi/2)\le m\). Integrating that derivative between the two endpoint \(\Phi\) values and using AA35 proves the assertion. The angle-profile inequality AA24 remains a still finer pointwise estimate, retaining every individual original Gram ratio.

\subsubsection{7. The literal tensor theta boundary of the angle energy}\label{endpoint-all_degree_angle_review--the-literal-tensor-theta-boundary-of-the-angle-energy}

This section supplies the source primitive behind AA7. Keep the original one-leg complex \(C_+=[V\xrightarrow{\Theta}\mathscr B]\) in degrees zero and one, with \(D=-x\partial_x\), and the original functions satisfying

\[
\mathcal MF_h=g/h,\qquad
\phi_*=(4\pi^2x^4-6\pi x^2)e^{-\pi x^2},\qquad
h(D)F_h=\Theta\phi_*.
\tag{AA38}
\]

The last identity follows by Mellin transformation: both sides transform to the exact \(g\), and Mellin uniqueness on these source spaces identifies them. The original source intertwining \(D\Theta=\Theta D\) lets all polynomial differential operators below act on the complex. Its inclusion in the original two-leg complex sends \(\phi\) to \((\phi,0)\) in degree zero and is identity in degree one, so the same primitive maps into the actual tau-base complex with its retained two-leg signs.

In the coefficient ring \(\mathbb C[s_1,\ldots,s_k]\), the cyclic definition of \(\chi\) gives

\[
\chi(s_1+\cdots+s_k)\in(h(s_1),\ldots,h(s_k)).
\tag{AA39}
\]

Indeed the tensor algebra \(E_h^{\otimes k}\) is this polynomial ring modulo that ideal, and its unit is the cyclic vector used to define \(\chi\). Evaluating \(\chi\) on the sum operator and that unit is zero by its actual annihilator definition.

The ideal expression can be chosen by a fully specified finite division. Start with \(F_0=\chi(s_1+\cdots+s_k)\); for \(j=1,\ldots,k\), divide \(F_{j-1}\) by the literal monic \(h(s_j)\), viewing the other variables as coefficients:

\[
F_{j-1}=h(s_j)L_j+F_j,\qquad\deg_{s_j}F_j<\deg h.
\tag{AA40}
\]

Later divisions do not increase the already restricted degrees in earlier variables, because \(h(s_j)\) has complex constant coefficients in those variables. Thus the final \(F_k\) lies in the span of the tensor remainder monomials, which are a basis of the quotient algebra. Its image in that algebra is zero by AA39, so \(F_k=0\). Summing the exact division identities proves

\[
\chi(s_1+\cdots+s_k)=\sum_{j=1}^k h(s_j)L_j(s_1,\ldots,s_k).
\tag{AA41}
\]

Every original coefficient of \(h\) and \(\chi\) remains in these finite quotient polynomials. This construction does not replace \(\chi\) by a product of simple factors or discard any nilpotent multiplicity.

For a polynomial \(P\) in the original sum variable, define the degree-\(k-1\) tensor cochain

\[
\begin{aligned}
\mathcal K_\chi(P)=\sum_{j=1}^k(-1)^{j-1}
&L_j(D_1,\ldots,D_k)P(D_1+\cdots+D_k)\\
&\bigl(F_h\otimes\cdots\otimes F_h
\otimes\phi_*\otimes F_h\otimes\cdots\otimes F_h\bigr),
\end{aligned}
\tag{AA42}
\]

where the \(j\)-th position is \(\phi_*\) in the degree-zero space \(V\), and every other position is the displayed \(F_h\) in degree one. There are \(j-1\) degree-one factors before that position. The tensor differential therefore acts on it with sign \((-1)^{j-1}\); this cancels the explicit coefficient sign in AA42. It acts trivially on the other degree-one factors. By AA38, by commutation of the polynomial operators with the differential, and then by AA41, one obtains exactly

\[
d\mathcal K_\chi(P)
=P(D_1+\cdots+D_k)\chi(D_1+\cdots+D_k)F_h^{\otimes k}
=\mathcal V_{h,k}(\chi P).
\tag{AA43}
\]

This proves the primitive with its full tensor signs. In particular, for \(x\in E\), regard \(X_{ij}x\) from AA7 as its literal polynomial in \(S\). Then

\[
\mathcal V_{h,k}(R_ix-R_jx)
=d\mathcal K_\chi(X_{ij}x),\qquad
G_i-G_j=X_{ij}^*H_jX_{ij}.
\tag{AA44}
\]

Thus the angle energy is the Gram energy of this exact original relation, while its cohomology class is zero. The unchanged arithmetic observation still satisfies \(J^{(k)}\mathcal V_{h,k}R_N=\eta_{h,k}\), with \(\eta_{h,k}[P]=\upsilon_h^{\otimes k}P(A_h^{(k)})1\) and the full unit \(\upsilon_h=j_h(g/h)\). Its further cohomology image is \(\sigma_h^{\otimes k}\eta_{h,k}\) in the original \(Q^{\otimes k}\). No source relation is dropped merely because those observations kill it.

At every nonempty original support label \(A\), lift the displayed coefficient maps by \((A,v)\mapsto(A,fv)\), and send external absence \(\tau\) to \(\tau\). The class of AA44 is the supported zero at the same label; its full source primitive and energy remain as written. This preserves the original cohomology over the tau base throughout the all-degree estimate.

\subsubsection{8. Full review of the root AGT1--AGT19 source}\label{endpoint-all_degree_angle_review--full-review-of-the-root-agt1agt19-source}

The complete corrected \nolinkurl{tau_all_degree_angle_transport_20260913.tex} was read at SHA-256 \nolinkurl{cc0c0c008b0532c5f120d1cea76b94e2d08153c87a694ce18775a5189c72d5f8}. It is accepted without a further mathematical correction. The initial source-cutoff sentence now states \(q-1\le N\le2q\), making its inclusion \(\mathcal P_N\subseteq H=\mathcal P_{2q}\) correct on the entire declared range.

The full acceptance covers the following calculations.

\begin{enumerate}
\def\labelenumi{\arabic{enumi}.}
\tightlist
\item
  AGT1--AGT4 retain the actual packet, full Taylor unit, original source and cohomology observations, both full masses, original vertical line, positive finite Grams, and all four signed determinant terms. The bounded source codomain is now exact.
\item
  AGT5--AGT8 agree with the independent derivation AA6--AA16. In the root source the cross matrix is written in the opposite orientation, \(U_j^*MU_i=G_j^{1/2}G_i^{-1/2}\); its adjoint-times-itself is precisely \(G_i^{-1/2}G_jG_i^{-1/2}\). Thus both presentations give the same squared singular values without commuting the square-root factors. The displayed two-dimensional projection block, all unit-cosine common vectors, and the remaining zero spectrum are correct. The primitive behind its relation-valued difference is supplied explicitly by AA38--AA44 above.
\item
  AGT9--AGT11 retain both full angle lists and remove only one exactly proved zero scalar contribution. The common ambient space has dimension \(2q-1\), including \(q=1\). The trace bound is the sum of the paired positive and negative Rayleigh bounds, so it needs no mutual commutation and introduces no extra factor two.
\item
  AGT12--AGT13 use the precise involution with variable \(k-S\) and conjugated coefficients. Its leading sign \((-1)^q\), its square, and its equality to complex conjugation on the original line are correct. The resulting integral is the full positive norm of the nonzero original polynomial \(\chi\), proving strict positivity on the actual moment path.
\item
  AGT14--AGT19 have the exact constant \(m=2q-1\). The derivative and tangent-line proof include the remaining zero-angle values by continuity and the actual positive mean \(\mathcal B/m\). The factor \(1/m\) in the derivative of \(\Phi_m\) cancels the factor \(m\) in the angle estimate. The integrated width equals the full original \(\log\kappa\), including repeated eigenvalues and endpoint eigenvalues equal to one. Both inverse hyperbolic branches and the explicit endpoint interval are correct. The source states the valid all-degree application, including actual quartet quotients, and preserves the absence of a new uniform estimate for \(\kappa\).
\end{enumerate}

The separate complete PA1--PA29 derivation in \nolinkurl{tau_principal_angle_volume_independent_review_20260913.md}, SHA-256 \nolinkurl{5f6a75afdbdfd4dcd81781a336702ea6d6545fd8fd0a9b2e014a605e414b3198}, was read in full. Its principal bases, projection spectra, positivity argument, and integrated constant agree with AA1--AA37. This provides an additional independent derivation, with its source scope retained; no finite numerical test is being substituted for these proofs.

\subsubsection{9. Reading, source pins, and mathematical scope}\label{endpoint-all_degree_angle_review--reading-source-pins-and-mathematical-scope}

The complete sealed AT source \nolinkurl{tau_arithmetic_determinant_transport_20260913.tex} at SHA-256 \nolinkurl{24cf7ce5be6a43cc656641b516185df10b2e4d7824600ece657120d1c549deae} and the complete AW source \nolinkurl{tau_relation_window_spectral_transport_20260913.tex} at SHA-256 \nolinkurl{c79e76c087efd59871e3234ff44fb18964515dbd34f0318dffba6b3f1440f6ec} were read in the preceding independent calculation. Their original projection and moment conventions are retained here. The full original \nolinkurl{arithmetic_input.tex} and \nolinkurl{Tau_Base_Cohomology/NOTE.md} were read earlier in this same audit and supply the stated theta complex, finite-packet unit, and tau-base tensor realization. AA6--AA44 provide the new finite angle, analytic transport, and explicit primitive proofs.

The companion receipt pins the final sources and the complete root-source acceptance above. These are closed written proofs and do not depend on a forthcoming supplement. No numerical or Lean execution, source normalization, remote mutation, or unproved uniform eigenvalue bound is claimed. The all-degree conclusion applies to the actual original moment Grams and includes every admissible quotient degree \(q\ge1\).

\endgroup


\endgroup
